% Generated by tools/build_new_dists_anthology_sections.py; do not edit by hand.

\section{Code-monomial matrix groups}\label{proof:new-dist:code_monomial}

\resultbox{\textbf{Canonical testing artifacts.}\par\smallskip Exact backend: \path{lambda_distributions/dists/code_monomial_sigma_mgfs/verification.py}.\par Regression: \path{lambda_distributions/dists/code_monomial_sigma_mgfs/test_verification.py}.\par Marimo: \path{marimo_notebooks/code_monomial.py}.}

\subsection*{Scope and outcome}

Let $C\leq(\Z/r\Z)^N$ be an additive code and let a coordinate-permutation
group $P\leq S_N$ preserve $C$.  We construct the honest unitary monomial
group $G(C,P)=D_C\rtimes P\leq U(N)$ and prove two equivalent formulas for
its complete generalized Molien, or $\sigma$-MGF, series: a weighted-cycle
average over $C\times P$ and a root-filtered sum over dual codewords that are
constant on permutation cycles.  We also prove the exact coefficient-orbit
interpretation, the complete-weight-enumerator specialization for $P=1$, and
the resulting coefficientwise convergence criterion and orbit-growth
obstruction.  An accompanying marimo notebook constructs five groups of
dimensions $2$ through $4$ and orders $4,8,24,18,64$.  All $95$ coefficients
through total degree five agree exactly; three independent evaluations of the
full series agree to $4.45\times10^{-16}$.

\subsection{Representation and target quantity}

Fix an integer $r\geq2$ and write $R=\Z/r\Z$ and
$\zeta=e^{2\pi i/r}$.  Let $C\leq R^N$ be an additive subgroup.  For
$c=(c_1,\ldots,c_N)\in C$, define
\[
 D(c)=\operatorname{diag}(\zeta^{c_1},\ldots,\zeta^{c_N}),
 \qquad D_C=\{D(c):c\in C\}.
\]
Let $P\leq S_N$ preserve $C$ under coordinate permutation, and use the
convention $P_\pi e_j=e_{\pi(j)}$.  Then conjugation by $P_\pi$ permutes the
diagonal entries of $D(c)$ and preserves $D_C$.  Hence
\[
 G(C,P)=\{D(c)P_\pi:c\in C,\ \pi\in P\}=D_C\rtimes P\leq U(N).
\]
The parametrization by $C\times P$ is injective, so $|G(C,P)|=|C||P|$.

For any finite $G\leq U(V)$ define
\begin{equation}\label{nd:code_monomial:eq:target}
 \MGF_{G,V}(\mathbf t)
 =\frac1{|G|}\sum_{g\in G}\prod_{i\geq1}\det(I-t_i g)^{-1}.
\end{equation}
For a partition $\tau=(\tau_1,\ldots,\tau_s)$, our symmetric-function
normalization is
\begin{equation}\label{nd:code_monomial:eq:coefficient-target}
 [m_\tau]\MGF_{G,V}
 =\dim\left(\bigotimes_{\alpha=1}^s\Sym^{\tau_\alpha}V\right)^G.
\end{equation}
Indeed, $[t^d]\det(I-tg)^{-1}$ is the character of $\Sym^dV$ at $g$;
averaging the product of these characters is the trace, hence the rank, of
the Reynolds projector onto the displayed invariant space.

\subsection{Weighted-cycle and equivariant dual-code formulas}

For $a\in R$, define the formal series
\begin{equation}\label{nd:code_monomial:eq:H-def}
 H_a(\mathbf t)=
 \sum_{\substack{\nu_1,\nu_2,\ldots\geq0\\
                  \sum_i\nu_i\equiv a\pmod r}}
 \prod_i t_i^{\nu_i}.
\end{equation}
Character orthogonality on $R$ gives the root filter
\begin{equation}\label{nd:code_monomial:eq:H-filter}
 \boxed{H_a(\mathbf t)=\frac1r\sum_{s=0}^{r-1}
 \zeta^{-as}\prod_i(1-\zeta^s t_i)^{-1}.}
\end{equation}
For a positive integer $\ell$, put
$H_a^{[\ell]}(\mathbf t)=H_a(t_1^\ell,t_2^\ell,\ldots)$, and define
\[
 C^\perp=\{u\in R^N:\langle c,u\rangle=0\text{ in }R
                         \text{ for every }c\in C\}.
\]

\begin{theorem}[Equivariant dual-code $\sigma$-MGF]\label{nd:code_monomial:thm:main}
For $G(C,P)$ on $V=\C^N$,
\begin{align}
 \MGF_{C,P}(\mathbf t)
 &=\frac1{|C||P|}\sum_{\pi\in P}\sum_{c\in C}
   \prod_{i\geq1}\prod_{Q\in\Cyc(\pi)}
   \left(1-\zeta^{\sum_{j\in Q}c_j}t_i^{|Q|}\right)^{-1},
   \label{nd:code_monomial:eq:weighted-cycle}\\
 &=\frac1{|P|}\sum_{\pi\in P}
   \sum_{\substack{u\in C^\perp\\
                    u\text{ constant on every cycle of }\pi}}
   \prod_{Q\in\Cyc(\pi)}H_{u_Q}^{[|Q|]}(\mathbf t),
   \label{nd:code_monomial:eq:dual-cycle}
\end{align}
where $u_Q$ is the common value of $u$ on $Q$.
\end{theorem}

\begin{proof}
Consider a cycle $Q=(j_1\,j_2\,\cdots\,j_\ell)$ of $\pi$.  On its
coordinate subspace, $D(c)P_\pi$ is a weighted cyclic shift.  Its $\ell$th
power is scalar multiplication by
$\zeta^{c_{j_1}+\cdots+c_{j_\ell}}$, and therefore
\[
 \det(I-tD(c)P_\pi\mid\C^Q)
 =1-\zeta^{\sum_{j\in Q}c_j}t^\ell.
\]
Multiplication over the disjoint cycles, over the variables $t_i$, and then
averaging over $C\times P$ proves \eqref{nd:code_monomial:eq:weighted-cycle}.

For a fixed cycle $Q$, expand its factor as
\[
 \prod_i\left(1-\zeta^{\sum_{j\in Q}c_j}t_i^{|Q|}\right)^{-1}
 =\sum_{\nu_1,\nu_2,\ldots\geq0}
  \zeta^{(\sum_i\nu_i)(\sum_{j\in Q}c_j)}
  \prod_i t_i^{|Q|\nu_i}.
\]
Assign to every $j\in Q$ the residue
$u_j=\sum_i\nu_i\pmod r$.  The resulting word $u$ is constant on every
cycle of $\pi$, and the total phase is $\zeta^{\langle c,u\rangle}$.
Orthogonality on the finite additive group $C$ gives
\[
 \frac1{|C|}\sum_{c\in C}\zeta^{\langle c,u\rangle}
 =\begin{cases}1,&u\in C^\perp,\\0,&u\notin C^\perp.\end{cases}
\]
After this filter, the sum of the allowed exponents on $Q$ is precisely
$H_{u_Q}^{[|Q|]}$.  Multiplying over cycles proves
\eqref{nd:code_monomial:eq:dual-cycle}.  This orthogonality argument is valid over
$\Z/r\Z$ for composite as well as prime $r$.
\end{proof}

\subsection{Pure diagonal groups and complete weight enumerators}

When $P=1$, every coordinate is a one-cycle.  Define the complete weight
enumerator of a code $A\leq R^N$ by
\[
 \cwe_A(X_0,\ldots,X_{r-1})
 =\sum_{u\in A}\prod_{j=1}^N X_{u_j}.
\]

\begin{corollary}[Pure diagonal specialization]\label{nd:code_monomial:cor:cwe}
For the diagonal group $D_C$,
\begin{equation}\label{nd:code_monomial:eq:cwe}
 \boxed{\MGF_{D_C}=\cwe_{C^\perp}(H_0,\ldots,H_{r-1}).}
\end{equation}
For $r=2$ this becomes
\[
 H_0=\tfrac12\left[\prod_i(1-t_i)^{-1}+\prod_i(1+t_i)^{-1}\right],
 \quad
 H_1=\tfrac12\left[\prod_i(1-t_i)^{-1}-\prod_i(1+t_i)^{-1}\right],
\]
so $\MGF_{D_C}=W_{C^\perp}(H_0,H_1)$ for the ordinary binary weight
enumerator $W$.
\end{corollary}

\begin{proof}
Set $P=1$ in \eqref{nd:code_monomial:eq:dual-cycle}.  The contribution of a fixed
$u\in C^\perp$ is $\prod_jH_{u_j}$, exactly the corresponding monomial in
the complete weight enumerator.
\end{proof}

This is the first verification family below.  It isolates the diagonal
character filter before any coordinate-orbit quotient is imposed.

\subsection{Exact coefficient formula}

For $\tau=(\tau_1,\ldots,\tau_s)$, let $\mathcal E_\tau(C^\perp)$ be the
set of matrices $A=(a_{\alpha j})\in\Z_{\geq0}^{s\times N}$ satisfying
\[
 \sum_{j=1}^N a_{\alpha j}=\tau_\alpha\quad(1\leq\alpha\leq s),
 \qquad
 \left(\sum_\alpha a_{\alpha1},\ldots,
       \sum_\alpha a_{\alpha N}\right)\bmod r\in C^\perp.
\]
The group $P$ acts on these arrays by permuting columns.

\begin{theorem}[Dual-code configuration orbits]\label{nd:code_monomial:thm:orbits}
For every partition $\tau$,
\begin{equation}\label{nd:code_monomial:eq:orbit-count}
 \boxed{[m_\tau]\MGF_{C,P}
 =\#\bigl(P\backslash\mathcal E_\tau(C^\perp)\bigr).}
\end{equation}
\end{theorem}

\begin{proof}
A monomial basis of
$\bigotimes_\alpha\Sym^{\tau_\alpha}\C^N$ is indexed by the exponent
arrays with the stated row sums.  The matrix $D(c)$ multiplies the basis
vector indexed by $A$ by
\[
 \zeta^{\sum_jc_j\sum_\alpha a_{\alpha j}}.
\]
It is fixed by every $D(c)$ exactly when the vector of total coordinate
exponents lies in $C^\perp$ modulo $r$.  Thus the $D_C$-invariant monomial
basis is indexed by $\mathcal E_\tau(C^\perp)$.  The quotient group $P$
permutes this basis by permuting columns, and its orbit sums form a basis of
the invariant space.  Equation \eqref{nd:code_monomial:eq:coefficient-target} completes the
proof.
\end{proof}

\subsection{Coefficientwise asymptotics}

Consider a sequence $(C_N,P_N)$ with a fixed modulus $r$.  The exact orbit
formula immediately gives
\begin{equation}\label{nd:code_monomial:eq:criterion}
 \MGF_{C_N,P_N}\text{ converges coefficientwise}
 \quad\Longleftrightarrow\quad
 \#\bigl(P_N\backslash\mathcal E_\tau(C_N^\perp)\bigr)
 \text{ is eventually constant for every fixed }\tau.
\end{equation}
Here ordinary convergence of a nonnegative integer sequence is equivalent to
eventual constancy.

There is a code-independent obstruction.  In each copy of $\Sym^r\C^N$,
the vectors $x_j^r$ are fixed by every $D(c)$ because $\zeta^{rc_j}=1$.
Consequently, for $\tau=(r,\ldots,r)$ with $s$ parts,
\begin{equation}\label{nd:code_monomial:eq:obstruction}
 \boxed{[m_{(r^s)}]\MGF_{C_N,P_N}
 \geq\#\bigl(P_N\backslash[N]^s\bigr).}
\end{equation}
Thus bounded orbit counts on every fixed Cartesian power of the coordinate
set are necessary for a raw coefficientwise limit.  Two immediate examples
are useful:
\begin{enumerate}
\item If $P_N=1$, then $[m_{(r)}]\MGF\geq N$, so every length-growing pure
diagonal family diverges in a fixed degree.
\item If $P_N$ is the regular cyclic shift group, then it has $N$ orbits on
ordered coordinate pairs, so $[m_{(r,r)}]\MGF\geq N$.
\end{enumerate}
Weight-enumerator convergence alone cannot control these coordinate-orbit
counts.  Conversely, \eqref{nd:code_monomial:eq:criterion} identifies the extra information
needed: bounded-degree \emph{equivariant} dual-code configurations.

\subsection{Independent verification strategy}

The executable verifier uses four logically separated calculations.
\begin{enumerate}
\item \textbf{Matrix construction.}  It creates every dense complex matrix
$D(c)P_\pi$, checks the expected group cardinality and distinctness, and
checks $g^*g=I$.
\item \textbf{Direct full-series value.}  At
$(t_1,t_2)=(0.07,0.11)$ it evaluates \eqref{nd:code_monomial:eq:target} using
\texttt{numpy.linalg.det} on the constructed matrices.
\item \textbf{Two full-series comparators.}  It evaluates
\eqref{nd:code_monomial:eq:weighted-cycle} from codeword cycle sums and
\eqref{nd:code_monomial:eq:dual-cycle} by enumerating $C^\perp$, testing cycle constancy, and
evaluating \eqref{nd:code_monomial:eq:H-filter}.
\item \textbf{Coefficient comparison.}  For the matrix route it obtains
$\operatorname{Tr}(\Sym^d g)$ from traces of actual matrix powers using
Newton's recurrence
\[
 d\,h_d(g)=\sum_{k=1}^d\operatorname{Tr}(g^k)h_{d-k}(g),\qquad h_0=1,
\]
then averages products over the matrix group.  The independent exact route
enumerates exponent arrays, imposes the dual-code condition, and counts
$P$-orbits as in Theorem~\ref{nd:code_monomial:thm:orbits}.
\end{enumerate}
The proof is general; these computations are regression tests, not a finite
substitute for the proof.

\subsection{Verification organized by group}

Every partition of every total degree from $0$ through $5$ was tested: $19$
coefficients for each of five groups, or $95$ exact comparisons in total.

\subsubsection{Pure diagonal group}

For $C=(\Z/2\Z)^2$ and $P=1$, the group has order $4$ on $\C^2$.  This
case separately tests Corollary~\ref{nd:code_monomial:cor:cwe}; in particular the degree-two
coefficient is $2$ before the coordinate swap identifies the two squares.

\subsubsection{Binary semidirect groups}

The full binary code of length two with $P=S_2$ produces the order-$8$
signed permutation group $W(B_2)$.  The even binary code of length three with
$P=S_3$ produces the order-$24$ group $W(D_3)$.  These test nontrivial cycle
types and dual codes of different dimensions.

\subsubsection{Ternary semidirect group}

For the ternary repetition code
$C=\{(a,a,a):a\in\Z/3\Z\}$ and $P=S_3$, the constructed group has order
$18$ on $\C^3$.  Its nonzero coefficients occur in a substantially different
degree pattern from the binary groups.

\subsubsection{Composite cyclic-alphabet group}

For
\[
 C=\{(a,b,a,b):a,b\in\Z/4\Z\}\leq(\Z/4\Z)^4
\]
and the cyclic coordinate group generated by $(1\,2\,3\,4)$, the constructed
group has order $64$ on $\C^4$.  This case directly tests the composite
modulus asserted in Theorem~\ref{nd:code_monomial:thm:main}.

\begin{center}
\begin{tabularx}{\linewidth}{@{}Xrrrrrc@{}}
\toprule
Group & $r$ & $N$ & $|C|$ & $|P|$ & $|G|$ & result\\
\midrule
$D_{(\Z/2)^2}$, pure diagonal &2&2&4&1&4&\textcolor{teal}{\textbf{PASS}}\\
$W(B_2)$ &2&2&4&2&8&\textcolor{teal}{\textbf{PASS}}\\
$W(D_3)$ &2&3&4&6&24&\textcolor{teal}{\textbf{PASS}}\\
ternary repetition $\rtimes S_3$ &3&3&3&6&18&\textcolor{teal}{\textbf{PASS}}\\
quaternary pair code $\rtimes C_4$ &4&4&16&4&64&\textcolor{teal}{\textbf{PASS}}\\
\bottomrule
\end{tabularx}
\end{center}

Selected coefficients are shown below; all remaining partitions through
degree five are checked by the notebook and test suite.
\begin{center}
\begin{tabularx}{\linewidth}{@{}Xrrrrrr@{}}
\toprule
Group & $(1)$ & $(2)$ & $(3)$ & $(2,1)$ & $(1,1,1)$ & $(4)$\\
\midrule
$D_{(\Z/2)^2}$ &0&2&0&0&0&3\\
$W(B_2)$ &0&1&0&0&0&2\\
$W(D_3)$ &0&1&1&1&1&2\\
ternary repetition $\rtimes S_3$ &0&0&3&4&5&0\\
quaternary pair $\rtimes C_4$ &0&0&0&0&0&3\\
\bottomrule
\end{tabularx}
\end{center}

At the full-series test point the direct matrix values are respectively
\[
 1.0507589401246002,\quad
 1.0254943428808700,\quad
 1.0286383823420668,\quad
 1.0106219737316153,\quad
 1.0013332933694528.
\]
Both formula routes reproduce every value.  Across all groups, the maximum
full-series discrepancy is $4.45\times10^{-16}$, the maximum coefficient
rounding residual is $2.23\times10^{-14}$, and the maximum unitarity residual
is $2.23\times10^{-16}$.

\subsection{Scope and reproduction}

The literal scalar model uses the cyclic additive alphabet $\Z/r\Z$.  For a
noncyclic alphabet such as the additive group of $\mathbf F_{p^s}$, there is
no faithful map of the entire alphabet into a single cyclic root group by
$a\mapsto\zeta^a$.  One must instead choose the appropriate family of
additive characters, character module, or Clifford--Weil representation.
Likewise, a monomial automorphism involving coordinate scalings, or a
semilinear automorphism involving a field automorphism, is not merely a
coordinate permutation in $P$ and needs a correspondingly generalized
cycle calculation.

It is also useful to separate three levels of code information.  For a pure
diagonal binary group, the ordinary Hamming weight enumerator of $C^\perp$
determines the answer through Corollary~\ref{nd:code_monomial:cor:cwe}.  For $r>2$, the
ordinary weight enumerator forgets the symbols and must be replaced by the
complete weight enumerator.  Once $P$ is nontrivial, even the complete weight
enumerator forgets how the coordinate group acts.  The exact required datum
is the equivariant collection of orbit counts in
\eqref{nd:code_monomial:eq:orbit-count}, equivalently a cycle-indexed complete-weight
enumerator.  Thus ordinary weight-enumerator convergence does not by itself
imply $\sigma$-MGF convergence.

The finite audit deliberately exercises the following boundaries:
\begin{enumerate}
\item $P=1$ versus nontrivial $P$, so the complete-weight and equivariant
formulas are tested separately;
\item prime moduli $2,3$ versus the composite modulus $4$;
\item full, even, repetition, and two-generator pair codes;
\item identity, symmetric, and cyclic coordinate groups.
\end{enumerate}
It does not claim to test every asymptotic code family mentioned in broader
surveys.  Equations \eqref{nd:code_monomial:eq:criterion} and \eqref{nd:code_monomial:eq:obstruction} are the
proved reusable tests for such a family once its codes and coordinate groups
are specified.

From the repository root, run
\begin{verbatim}
.venv/bin/python -m lambda_distributions.dists.code_monomial_sigma_mgfs.verification
.venv/bin/pytest -q lambda_distributions/dists/code_monomial_sigma_mgfs/test_verification.py
.venv/bin/marimo edit marimo_notebooks/code_monomial.py
\end{verbatim}
The source and PDF in this directory are standalone artifacts intended for
later inclusion.  No combined PDF is edited or rebuilt.
